\documentclass[11pt]{article}
% Shared local declarations. Document roots use the repository's business-admin family.
\usepackage{business-admin}
\usepackage{amsmath,amssymb,tabularx,longtable,array,booktabs,enumitem}
\usepackage{ragged2e}
\newcolumntype{Y}{>{\RaggedRight\arraybackslash}X}
\newcommand{\SuiteMeta}[2]{\def\DocTitle{#1}\def\DocSubtitle{#2}\def\DocVersion{1.0}\def\DocStatus{Educational and operational reference}\def\DocOwner{Parliamentary Procedure and Governance Suite}\def\DocSummary{Independently authored reference and working materials, based on RONR 12th edition.}}
\newcommand{\SuiteOpen}{\makedoctitle\vspace{-0.5em}\noindent\rule{\linewidth}{0.4pt}\par\medskip}
\newcommand{\RonrAnchor}[1]{\par\smallskip{\footnotesize\textbf{Primary reading:} Robert's Rules of Order Newly Revised, 12th ed., #1. Original synthesis; not official interpretation.}\par}
\newcommand{\WriteLine}{\par\noindent\rule{\linewidth}{0.35pt}\par\vspace{0.5cm}}
\newcommand{\WriteLines}[1]{\par\noindent\rule{\linewidth}{0.35pt}\par\vspace{0.35cm}\noindent\rule{\linewidth}{0.35pt}\par\vspace{0.35cm}\noindent\rule{\linewidth}{0.35pt}\par\vspace{0.25cm}}
\newcommand{\Field}[1]{\par\medskip\noindent\textbf{#1}\par\WriteLines{3}}
\newcommand{\GovCheckbox}{\(\square\)}
\setlength{\emergencystretch}{3em}

\SuiteMeta{Motion Classification and Order of Precedence — Template}{Guided exercise | fictional example}
\begin{document}
\SuiteOpen
\section{Method}
Example: previous question is a subsidiary motion used to close debate; a two-thirds vote is the ordinary requirement. Consult Section 16 for exact conditions.

Work step by step. Use separate sheets for each procedural branch and verify each choice against the applicable section of the twelfth edition. The goal is a reproducible reasoning trail, not merely a final yes/no answer.

\section{Practice sequence}
\begin{enumerate}[leftmargin=*,itemsep=4pt]
\item Identify the fact pattern and meeting’s controlling authority.
\item Determine current procedural state and any exception.
\item Apply the relevant section and compute any decision threshold.
\item Write the action or ruling in plain procedural language.
\item Use the blank companion to record an independent attempt.
\end{enumerate}

\section{Fictional worked starting point}
Example: previous question is a subsidiary motion used to close debate; a two-thirds vote is the ordinary requirement. Consult Section 16 for exact conditions.

\section{Self-check}
\begin{itemize}[leftmargin=*,itemsep=4pt]
\item Did I separate a member’s right from the assembly’s power?
\item Did I use the right population and threshold?
\item Does the record contain the exact final action?
\item Could another parliamentarian reconstruct the sequence?
\end{itemize}
\end{document}
